%% LyX 2.4.4 created this file.  For more info, see https://www.lyx.org/.
%% Do not edit unless you really know what you are doing.
\documentclass[english,handout]{beamer}
\usepackage{mathptmx}
\usepackage[T1]{fontenc}
\usepackage[latin9]{inputenc}
\usepackage{units}
\usepackage{amssymb}
\usepackage{graphicx}

\makeatletter
%%%%%%%%%%%%%%%%%%%%%%%%%%%%%% Textclass specific LaTeX commands.
% this default might be overridden by plain title style
\newcommand\makebeamertitle{\frame{\maketitle}}%
% (ERT) argument for the TOC
\AtBeginDocument{%
  \let\origtableofcontents=\tableofcontents
  \def\tableofcontents{\@ifnextchar[{\origtableofcontents}{\gobbletableofcontents}}
  \def\gobbletableofcontents#1{\origtableofcontents}
}

%%%%%%%%%%%%%%%%%%%%%%%%%%%%%% User specified LaTeX commands.
\usetheme{Warsaw}
% or ...

\setbeamercovered{transparent}
% or whatever (possibly just delete it)
\setbeamercovered{invisible} %makes overlays invisible


%gets rid of navigation symbols
\setbeamertemplate{navigation symbols}{}


\usepackage{pgfpages}
\pgfpagesuselayout{4 on 1}[a4paper, landscape, border shrink=7mm]
\pgfpageslogicalpageoptions{1}{border code=\pgfusepath{stroke}}
\pgfpageslogicalpageoptions{2}{border code=\pgfusepath{stroke}}
\pgfpageslogicalpageoptions{3}{border code=\pgfusepath{stroke}}
\pgfpageslogicalpageoptions{4}{border code=\pgfusepath{stroke}}

\makeatother

\usepackage{babel}
\begin{document}
\title[Properties of Power Series]{Ch. 10.2: Properties of Power Series}
\subtitle{Calculus III | MTH 331~~|~~Section A~~|~~Fall 2012}
\author{Dr.\,James Ashe}
\titlegraphic{\includegraphics[height=1.5cm]{../../jcsu_bull.jpg}}
\institute{Johnson C. Smith University}
\date{{}}

\makebeamertitle

%\beamerdefaultoverlayspecification{<+->}
\begin{frame}{}

\begin{block}{Definition}

A \emph{power series }is an infinite series of the form \\
\hfill${\displaystyle \sum_{k=0}^{\infty}c_{k}x^{k}=c_{0}+c_{1}x+c_{2}x^{2}+\dots+c_{n}x^{n}+c_{n+1}x^{n+1}\dots}$\hfill\,\\
or, more generally,\\
\hfill${\displaystyle \sum_{k=0}^{\infty}c_{k}(x-a)^{k}=c_{0}+c_{1}(x-a)+\dots+c_{n}(x-a)^{n}+\dots}$\hfill\,\\
where $c_{k}$ are constants.
\end{block}
\begin{itemize}
\item <2->The constant $a$ is the \textbf{center}.
\item <3->The set of $x$'s such that the series converges is called the
\textbf{interval of convergence}, $I=[c,d]$. i.e., \\
$I=\{x\in\mathbb{R}|\,\sum_{k=0}^{\infty}c_{k}(x-a)^{k}<\infty\}$.
\item <4->The \textbf{radius of convergence }of the series, $R$, is the
distance from $a$ to $d$, i.e., $R=|d-a|$.
\end{itemize}
\end{frame}

\begin{frame}{Finding an interval of convergence}

\begin{itemize}
\item [ex.]<2->${\displaystyle \sum\frac{x^{k}}{2^{k}}}$ 
\item []<3->Need to find all $x$ such that the series converges. \uncover<4->{By
the Ratio Test, ${\displaystyle \lim_{k\rightarrow\infty}\left|\frac{a_{n+1}}{a_{n}}\right|}={\displaystyle \lim_{k\rightarrow\infty}\left|\frac{\nicefrac{x^{k+1}}{2^{k+1}}}{\nicefrac{x^{k}}{2^{k}}}\right|}={\displaystyle \lim_{k\rightarrow\infty}\frac{\left|x\right|}{2}=\frac{\left|x\right|}{2}}$
}
\item []<4->So this series diverges when$\frac{\left|x\right|}{2}>1\Rightarrow x<-2\mbox{ and }x>2$.
\item <5->Need to test the endpoints. 
\item []<6->At $x=-2$, ${\displaystyle \sum_{k=0}^{\infty}\frac{(-2)^{k}}{2^{k}}=1-1+1-1+\cdots}$,
so it \uncover<7->{diverges.}
\item []<8->At $x=2$, $\sum_{k=0}^{\infty}\frac{(2)^{k}}{2^{k}}=1+1+\cdots$,
so it \uncover<9->{diverges.}
\item []<10->So the interval of convergence is \uncover<11->{$-2<x<2$.}
\end{itemize}
\end{frame}

\begin{frame}{}

\begin{itemize}
\item [ex.]${\displaystyle \sum_{k=1}^{\infty}\frac{(x-2)^{k}}{nk^{n}}}$
\item []<2->Using the Ratio Test, ${\displaystyle \lim_{k\rightarrow\infty}\left(\frac{n}{n+1}\cdot\frac{5^{n}}{5^{n+1}}\cdot\frac{\left|x-2\right|^{n+1}}{\left|x-2\right|^{n}}\right)=\frac{1}{5}\left|x-2\right|}$.
\item []<3->$\frac{1}{5}\left|x-2\right|>1\Rightarrow\left|x-2\right|>5$,
so then the series diverges when $x\not\in\left(-3,7\right)$.
\item []<4->For $x=7$, the series is ${\displaystyle \sum_{k=1}^{\infty}\frac{1}{k}}$
is harmonic so it diverges.
\item []<5->For $x=-3$, the seres is $\sum_{k=1}^{\infty}(-1)^{k}\frac{1}{k}$,
which is the alternating harmonic series which converges (by the Alternating
Series Test). 
\item []<6->So the interval of convergence is $-3\leq x<7$.
\end{itemize}
\end{frame}

\begin{frame}{Finding a power series representation for $f(x)$ and its interval
of convergence}

\begin{exampleblock}{}

\begin{center}$g(x)=\frac{1}{1-x}={\displaystyle \sum_{k=0}^{\infty}x^{k}}$
when $|x|<1$.\end{center}
\end{exampleblock}
\begin{itemize}
\item The geometric series, ${\displaystyle \sum_{k=0}^{\infty}x^{k}}=\frac{1}{1-x}$
is a power series (replacing the $r$ with a variable; each $c_{i}=1$).
\item [ex.]<2-> Find a power series for $f(x)=\frac{1}{1+x^{3}}$ centered
at $0$.
\item []<3->Substitute $-x^{3}$ for $x$ in $g(x)$. This gives $\frac{1}{1-(-x^{3})}$,
and
\item []<4->$\frac{1}{1-(-x^{3})}={\displaystyle \sum_{k=0}^{\infty}(-x^{3})^{n}}$
\emph{provided} that $|-x^{3}|<1.$
\item []<5->Now \textrm{${\displaystyle \sum_{k=0}^{\infty}(-x^{3})^{n}=\sum_{k=0}^{\infty}(-1)^{n}x^{3n}}$
and $|x^{3}|<1\Rightarrow|x|<1$.}
\end{itemize}
\end{frame}

\begin{frame}{}

\begin{itemize}
\item [ex. 33]Find the power series for $h(x)=\frac{1}{(1-x)^{2}}$ using
$g(x)=\frac{1}{1-x}$.
\item []<2->Note: $\frac{1}{(1-x)^{2}}=\frac{d}{dx}\left(\frac{1}{1-x}\right)$,
and 
\item []<3->$\frac{1}{1-x}={\displaystyle \sum_{k=0}^{\infty}x^{k}}$ when
$\left|x\right|<1$.
\item []<4->So then $h(x)=\frac{1}{(1-x)^{2}}=\frac{d}{dx}\left(\frac{1}{1-x}\right)=\frac{d}{dx}{\displaystyle \sum_{k=0}^{\infty}x^{k}}$
\uncover<5->{$={\displaystyle \sum_{k=1}^{\infty}kx^{k-1}}$}. 
\end{itemize}
\end{frame}

\begin{frame}{}

\textbf{Homework }due (10/24)

\textbf{10.2: }9-12, 21-24, 34-35, 39-41

\end{frame}

\end{document}
